% !PRINTED 140
\weberpage{140}
%<<CONT

nach Potenzen von $x$ ordnen, so sind die Coëfficienten der einzelnen Potenzen von $x$ gleichfalls symmetrische Functionen. Denn sie ändern sich bei der Vertauschung der $\alpha$ ebenso wenig wie die Function $f(x)$ (\S.\,1). Wir setzen
\begin{equation}
 f(x)=x^n+a_1x^{n-1}+a_2x^{n-2}+\cdots+a_{n-1}x+a_n.
\tag{2}
\end{equation}

Die Functionen $a_1,a_2,\ldots,a_n$ sind die Coëfficienten der algebraischen Gleichung, deren Wurzeln $\alpha_1,\alpha_2,\ldots,\alpha_n$ sind. Sie werden die \gesperrt{symmetrischen Grundfunctionen} genannt. Diese symmetrischen Grundfunctionen haben wir schon im \S.\,7 gebildet; es ist, wie wir uns erinnern, $a_1$ die mit negativem Zeichen genommene Summe der $\alpha$, $a_2$ die Summe der Producte zu zweien, $a_3$ die mit negativem Zeichen genommene Summe der Producte zu dreien u. s. f., endlich $a_n$ das Product sämmtlicher $\alpha$, mit positivem oder negativem Zeichen, je nachdem $x$ gerade oder ungerade ist.
% [sic] The registered source prints x here; the parity determining the sign is that of the degree n.
Wir setzen abkürzend
\begin{equation}
\begin{aligned}
 a_1&=-\sum \alpha_1,\\
 a_2&=+\sum \alpha_1\alpha_2,\\
 a_3&=-\sum \alpha_1\alpha_2\alpha_3,\\
 &\ \vdots\\
 a_n&=\pm \alpha_1\alpha_2\cdots\alpha_n.
\end{aligned}
\tag{3}
\end{equation}

Das Ziel unserer Betrachtungen ist der Beweis des Hauptsatzes, \gesperrt{dass alle symmetrischen Functionen der $\alpha$ sich rational durch die Grundfunctionen ausdrücken lassen.}

\begin{center}
{\large \S.\ 42.}

\vspace{0.5em}
{\large\gesperrt{Die Potenzsummen.}}
\end{center}

Wir beschäftigen uns zunächst mit einer anderen speciellen Art symmetrischer Functionen, den \gesperrt{Potenzsummen}. Bedeutet nämlich $\nu$ irgend einen ganzzahligen positiven Exponenten, so gehört
\begin{equation}
 s_\nu=\alpha_1^\nu+\alpha_2^\nu+\cdots+\alpha_n^\nu
\tag{1}
\end{equation}
offenbar zu den symmetrischen Functionen, und $s_\nu$ wird die $\nu$te Potenzsumme genannt. Wir wollen Formeln ableiten, nach denen die Potenzsummen durch die Grundfunctionen ausdrückbar sind.

Wir bezeichnen in der Folge immer, wenn $\varphi(x)$ irgend eine Function von $x$ ist, mit
\[
 S[\varphi(\alpha)]
\]
%CONT>>
